\section{Implementación del Sistema Odoo 19}

La implementación del sistema Odoo 19 en \textbf{Machu Picchu Exclusive Tours E.I.R.L.} se realizó bajo la modalidad \textbf{Odoo Online (SaaS)}, lo que permitió desplegar una solución centralizada en la nube capaz de soportar la operación turística sin requerir infraestructura de servidores locales ni costos elevados de mantenimiento TI.

Esta modalidad garantiza la autonomía de los datos y la escalabilidad del sistema, permitiendo a la empresa gestionar su crecimiento en número de clientes y reservas sin cambios estructurales en la plataforma. Asimismo, asegura el acceso remoto seguro mediante credenciales únicas para cada perfil de usuario (\textbf{Dirección, Comercial, Operaciones, Administración}), facilitando la supervisión desde cualquier ubicación.

El sistema fue configurado estrictamente de acuerdo con los flujos operativos de la agencia, priorizando tres pilares fundamentales:

\begin{enumerate}
    \item \textbf{Visión 360° del PAX:} Centralización de datos de contacto, documentación y notas en una ficha única de cliente.
    \item \textbf{Control del flujo de reservas:} Uso de tableros Kanban para seguir cada viaje desde la cotización hasta la ejecución.
    \item \textbf{Control de adelantos y saldos:} Registro ordenado de pagos y saldos antes de ejecutar compras no reembolsables.
\end{enumerate}

\subsection{Configuración inicial}

En esta etapa se realizó la creación del entorno de trabajo en \textbf{Odoo Online}, definiendo:

\begin{itemize}
    \item La base de datos principal de la empresa.
    \item La configuración inicial de idioma, moneda, zona horaria y parámetros generales.
    \item La creación de usuarios, perfiles de acceso y permisos según los roles identificados.
\end{itemize}

Se activaron los módulos definidos en el alcance (CRM, Proyectos, actividades y ventas/facturación básica), asegurando que el entorno se encontrara preparado para la parametrización de procesos y la carga de datos maestros (clientes y servicios turísticos).

\subsection{Implementación del Módulo CRM Turístico}

En el módulo CRM se configuraron:

\begin{itemize}
    \item Fichas de contacto adaptadas al contexto turístico, permitiendo registrar información del PAX y del viaje.
    \item Estructura para adjuntar de forma segura \textbf{fotos de pasaportes} y otros documentos relevantes.
    \item Registro de notas internas e historial de interacciones (llamadas, acuerdos, cambios de itinerario).
    \item Uso de etiquetas para segmentar contactos según mercados de origen, tipo de servicio o campañas.
\end{itemize}

\subsection{Implementación del Módulo de Proyectos}

Se implementó un tablero Kanban específico para el control de reservas, con las etapas:

\begin{enumerate}
    \item Borrador de Cotización.
    \item Confirmación (Adelanto).
    \item Compra de Tickets/Hoteles.
    \item Ejecución (Venta Final).
\end{enumerate}

Cada reserva se gestiona como un proyecto o tarea con:

\begin{itemize}
    \item Asignación de responsable interno.
    \item Documentos adjuntos (comprobantes, reservas, vouchers).
    \item Actividades asociadas (llamadas, correos, recordatorios) alineadas al itinerario.
\end{itemize}

\subsection{Gestión de Actividades y Alertas}

Se configuró el uso de la funcionalidad estándar de \textbf{Actividades} de Odoo para:

\begin{itemize}
    \item Programar recordatorios de llamadas y correos clave (confirmación de adelantos, envío de vouchers, seguimientos).
    \item Visualizar tareas pendientes en la \textbf{vista de calendario}, facilitando la planificación diaria.
    \item Asegurar el seguimiento oportuno antes de hitos críticos del viaje (inicio del tour, traslados, visitas principales).
\end{itemize}

\subsection{Gestión de Cobros y Saldos}

En el módulo de ventas/facturación básica se configuró:

\begin{itemize}
    \item Registro manual de \textbf{adelantos} (entre el 30\% y 50\% del valor del paquete) y saldos.
    \item Emisión de comprobantes internos o recibos para confirmar la \textbf{recepción de fondos} al turista.
    \item Visualización del estado de pago (Pagado / Saldo pendiente) asociado a cada reserva, como \textbf{condición previa} a la compra de boletos o servicios no reembolsables.
\end{itemize}

